\section{Document extraction and provenance}\label{section:disc-doc-extraction-and-provenance}

\textbf{Interpreting the rubric scores (RQ1).} The pinned configuration reaches a strict-F1 score of $0.838$ for tables, and $0.840$ for figures (Table~\ref{tab:table-emission-metrics}, Table~\ref{tab:figure-emission-metrics}). The sources of improvement are important, and should be interpreted carefully: off-the-shelf Docling already localizes tables effectively, with crop and mask F1 scores of $90.1\%$ and $95.8\%$, yet reaches only $0.366$ strict table-F1. The Docling-only baseline (v01), with its post-processing steps, marginally improves the strict-F1, to $0.400$. Both variants perform poorly on the strict rubric because of the same reason: they do not include table footnotes. The rise to $0.838$ strict-F1 is therefore attributable almost entirely to the footnote expansion, which expands table crops downward by a footnote-expansion multiplier of $\times1.2$. This correction raises footnote precision from $43.8\%$ to $91.4\%$, showing that the main bottleneck was not table localization, but footnote capture.

There is a similar story for figures. Off-the-shelf Docling emits many decorative icons, and labels them as figures. It also often fails to attach the correct caption to each figure. These two common failure modes cause the off-the-shelf Docling to reach only $0.447$ strict figure-F1. The production pipeline addresses these failures through post-processing: icon and size filtering, sub-figure merging, and geometric caption matching. The post-processing step increases the strict figure-F1 to $0.840$. When taken together, the results summarize the document-extraction strategy of the thesis: the competent off-the-shelf parser is not replaced, but rather, its well-understood failure modes are addressed, and better performance is achieved for downstream use.

\medskip{}

\textbf{The remaining errors are bounded and interpretable.} The remaining table failure modes fall into two main groups. The first consists of caption-attachment errors (table continuation markers parsed as separate tables, multi-page and rotated tables, captions attached to the side or corner of a table); the second consists of cropping-geometry errors, where boundaries of side-by-side tables cannot be correctly detected, and are falsely merged, or cropped incorrectly. 

The remaining figure failures are dominated by decorative icons that are falsely detected as figures, making up 14 of 21 false positives. It is important to note that these false positives are still masked; their text does not leak into the main text content, so downstream stages are not affected. Therefore, the remaining failure modes are understood as future work, and they are scoped: they are enumerable, and isolated from the knowledge-extraction stage. Document-extraction does not need to be perfect, it only needs to fail in ways that do not corrupt the downstream stages.

\medskip{}

\textbf{Provenance is guaranteed by construction (RQ2).} Provenance is a structural property of the pipeline; it is not a quality estimate. It holds completely across the full ingested corpus, consisting of $977$ documents, $35\,896$ text elements, $4\,479$ figures, and $1\,960$ tables. Every populated field can be traced back to its source. Every paragraph carries a \texttt{unique\_path} (PMCID, section path, position). Each figure and table retains its cropped image and bounding box ($100\%$). Every table carries a caption; for figures, this number drops to $86.3\%$, which is mostly caused by the small decorative icons that are wrongly emitted as figures. 

Provenance is built into the storage representation of the documents from the beginning (Subsection~\ref{subsec:doc-provenance-preserving-storage}). A pipeline that discards information like bounding box coordinates or source paths cannot recover them later. This document-level provenance supports the stronger claim made in the knowledge-extraction stage: every extracted rule is traceable back to the evidence it was produced from.